\documentclass[10pt,twocolumn]{article}

\usepackage[margin=1in]{geometry}
\usepackage{amsmath,amssymb}
\usepackage{graphicx}
\usepackage{booktabs}
\usepackage{colortbl}
\usepackage{hyperref}
\usepackage{xcolor}
\usepackage{microtype}
\usepackage{listings}
\usepackage{caption}
\usepackage{subcaption}
\usepackage{abstract}
\usepackage{titlesec}
\usepackage{parskip}
\usepackage{fancyhdr}
\usepackage{natbib}
\usepackage[section]{placeins}  % FloatBarrier at every \section

\definecolor{codegray}{rgb}{0.95,0.95,0.95}
\definecolor{codeblue}{rgb}{0.1,0.2,0.6}

\lstset{
  backgroundcolor=\color{codegray},
  basicstyle=\ttfamily\small,
  keywordstyle=\color{codeblue}\bfseries,
  language=Python,
  frame=single,
  framerule=0pt,
  breaklines=true,
  captionpos=b,
}

\hypersetup{
  colorlinks=true,
  linkcolor=codeblue,
  citecolor=codeblue,
  urlcolor=codeblue,
}

\pagestyle{fancy}
\fancyhf{}
\rhead{\small Adaptive Precision Attention}
\lhead{\small Chaithu Talasila}
\cfoot{\thepage}

\title{\textbf{Adaptive Precision Attention:\\
Entropy-Guided Quantization via Evolutionary Search}}

\author{
  Chaithu Talasila\\
  \texttt{https://github.com/themoddedcube/adaptive-precision-attention}
}

\date{}

\begin{document}

\maketitle
\thispagestyle{fancy}

%------------------------------------------------------------------
\begin{abstract}
We present an evolutionary approach to discovering per-block precision
allocation policies for mixed-precision attention. Using OpenEvolve, an
LLM-guided evolutionary code search framework, we evolve a precision
policy function over a search space of 14 per-block statistical features
and 6 numerical precision levels. After 80 iterations across a population
of 50 candidates, the search converges on a remarkably simple two-threshold
policy: use INT8 quantization for all attention blocks except those
exhibiting both statistical outliers and low softmax entropy (below 2.0),
which retain FP16 precision. Evaluated across 12 workloads spanning
sequence lengths from 512 to 8{,}192, the policy compresses 58\% of
attention blocks from 16-bit to 8-bit representation with no measurable
accuracy degradation, while a naive uniform INT8 baseline suffers
catastrophic failure ($>$1000$\times$ worse RMSE) on outlier-heavy
workloads. The simplicity of the converged policy---two comparisons and an
AND gate---makes it directly implementable as a hardware precision
controller for ASIC attention accelerators. We further demonstrate an
entropy-guided KV-cache quantization module achieving approximately
$2\times$ memory compression for non-outlier workloads.
\end{abstract}

%------------------------------------------------------------------
\section{Introduction}

The attention mechanism is the computational core of transformer-based
language models, and its efficient implementation has been a sustained
focus of systems research. FlashAttention \citep{dao2022flashattention}
and its successors \citep{dao2023flashattention2,shah2024flashattention3}
achieve near-peak hardware utilization through tiled, fused computation
that avoids materializing the full $N \times N$ attention matrix.
FlashAttention-3 introduces FP8 quantization on Hopper GPUs, achieving
up to 1.2 PFLOPs/s, but treats precision as a global choice: all blocks
are computed at the same numerical format.

In practice, different blocks of the attention matrix have different
precision requirements. Blocks where the softmax distribution is peaked
and input values contain statistical outliers are sensitive to
quantization error. Blocks where attention is uniformly distributed are
robust to aggressive compression. This motivates \emph{adaptive precision
attention}, where each block is quantized to the minimum precision that
preserves output quality.

The challenge is determining which blocks need protection. We pose
precision allocation as a program synthesis problem and solve it with
evolutionary search over per-block statistics.

\paragraph{Contributions.}
\begin{itemize}
  \item An evolutionary framework for precision policy discovery over 14
  per-block statistics and 6 precision levels.
  \item The entropy threshold finding: softmax entropy and outlier
  presence are the only features that matter; the gap is bimodal and wide.
  \item A hardware-friendly precision controller: two comparisons, one AND
  gate, suitable for direct ASIC implementation.
  \item An entropy-guided KV-cache module achieving $\approx 2\times$
  memory compression with $<$0.02 mean absolute reconstruction error.
\end{itemize}

%------------------------------------------------------------------
\section{Related Work}

\paragraph{Efficient attention.}
FlashAttention \citep{dao2022flashattention} introduced tiled, IO-aware
attention computation. FlashAttention-2 \citep{dao2023flashattention2}
improved parallelism. FlashAttention-3 \citep{shah2024flashattention3}
targets Hopper GPUs with warp specialization and uniform FP8 block
quantization. All three treat precision as a global setting.

\paragraph{Adaptive quantization.}
\citet{boroujeni2026adaptive} propose per-token KV-cache quantization
with four precision levels selected by entropy and attention variance,
achieving 17.75\% latency reduction. SmoothQuant
\citep{xiao2023smoothquant} redistributes quantization difficulty between
activations and weights. GPTQ \citep{frantar2023gptq} uses second-order
information for weight quantization. Our work focuses on attention
computation precision using evolutionary search, targeting ASIC
implementation.

\paragraph{Program synthesis.}
OpenEvolve \citep{openevolve2025} extends FunSearch
\citep{romera2024funsearch} with multi-objective optimization and
island-based population management. We use it as the search engine with
Claude Code as the LLM backbone.

%------------------------------------------------------------------
\section{Method}

\subsection{Problem Formulation}

We define a precision policy as a function mapping per-block statistics
to a precision level:
\begin{align*}
\text{policy}: \mathcal{S} \rightarrow \{&\texttt{fp16},\;\texttt{fp8\_e4m3},\;\texttt{fp8\_e5m2},\\
                                         &\texttt{int8},\;\texttt{fp4},\;\texttt{int4}\}
\end{align*}

For each query block $\mathbf{Q}_i$ and key-value block
$(\mathbf{K}_j, \mathbf{V}_j)$, we compute block statistics, query the
policy, and apply a quantize-dequantize roundtrip before accumulation
(always in FP32).

\subsection{Block Statistics}

For each block pair $(i,j)$ we extract 14 statistics: L2 norms of
$\mathbf{Q}$, $\mathbf{K}$, $\mathbf{V}$ blocks; statistics of
pre-softmax scores $\mathbf{S} = \mathbf{Q}\mathbf{K}^\top/\sqrt{d}$;
the Shannon entropy of the row-wise softmax; causal distance; block
indices; and an outlier flag:
\[
\texttt{has\_outlier} = \max|\mathbf{S}| > \mu_{|\mathbf{S}|} + 10\,\sigma_{|\mathbf{S}|}
\]

Block size is 128 tokens, matching FlashAttention-3's tile size.

\subsection{Evolutionary Search}

We use OpenEvolve \citep{openevolve2025} with 50 candidates across 3
islands, 80 iterations of diff-based mutation, Claude Sonnet (80\%) and
Haiku (20\%) as LLM backbone, and a fitness function weighting 50\%
accuracy ($1/(1 + \text{RMSE} \cdot 10)$ vs FP64) + 30\% compression +
20\% reliability.

\subsection{Evaluation Workloads}

We evaluate on 12 workloads spanning $N \in \{512, 2048, 4096, 8192\}$,
$d \in \{64, 128\}$, with and without causal masking, and with 0.1\% of
elements scaled by $100\times$ to simulate LLM activation outliers.

%------------------------------------------------------------------
\section{Results}

\subsection{Converged Policy}

After 80 iterations (combined score 0.745, target $>$0.65), the search
converged on:

\begin{lstlisting}[caption={Evolved precision policy (2 of 14 features used).}]
def precision_policy(block_stats):
    has_outlier = block_stats.get('has_outlier', False)
    entropy = block_stats.get('entropy', 4.3)
    if has_outlier and entropy < 2.0:
        return "fp16"
    return "int8"
\end{lstlisting}

All intermediate formats (FP8, FP4, INT4) were eliminated by the search.

\subsection{Accuracy vs.\ Baselines}

Table~\ref{tab:avg_results} summarises aggregate metrics.
Table~\ref{tab:results} gives per-workload detail.
Figure~\ref{fig:accuracy} shows the log-scale RMSE comparison.
On outlier workloads, uniform INT8 fails catastrophically (RMSE $>$ 1.0),
while the evolved policy matches FP16 exactly.

\begin{table}[tbp]
\centering
\caption{Average metrics across all 12 workloads.}
\label{tab:avg_results}
\small
\begin{tabular}{lrrc}
\toprule
Strategy & Avg.\ RMSE & Bits & Compression \\
\midrule
Uniform FP16   & 4.93e-4 & 16.0 & 1.0$\times$ \\
Evolved Policy & 1.08e-2 & 11.3 & 1.4$\times$ \\
Uniform INT8   & 5.61e-1 &  8.0 & 2.0$\times$ \\
\bottomrule
\end{tabular}
\end{table}

\begin{table*}[tbp]
\centering
\caption{Per-workload relative RMSE vs.\ FP64 reference.
  Shaded rows are outlier workloads. Uniform INT8 is catastrophic there;
  evolved policy matches FP16.}
\label{tab:results}
\small
\begin{tabular}{ccccrrrc}
\toprule
$N$ & $d$ & Causal & Outliers & FP16 RMSE & INT8 RMSE & Evolved RMSE & Bits \\
\midrule
  512 &  64 & No  & No  & 1.84e-4 & 1.81e-2 & 1.81e-2          &  8 \\
\rowcolor{gray!15}
  512 & 128 & Yes & Yes & 5.74e-4 & 1.33e+0 & \textbf{5.74e-4} & 16 \\
\rowcolor{gray!15}
 2048 &  64 & No  & Yes & 1.00e-3 & 1.35e+0 & \textbf{1.00e-3} & 16 \\
 2048 & 128 & Yes & No  & 1.67e-4 & 1.79e-2 & 1.79e-2          &  8 \\
 8192 &  64 & Yes & No  & 1.73e-4 & 1.62e-2 & 1.62e-2          &  8 \\
\rowcolor{gray!15}
 8192 & 128 & No  & Yes & 1.14e-3 & 1.38e+0 & \textbf{1.14e-3} & 16 \\
\rowcolor{gray!15}
 2048 &  64 & Yes & Yes & 1.06e-3 & 1.27e+0 & \textbf{1.06e-3} & 16 \\
 2048 & 128 & No  & No  & 1.82e-4 & 1.92e-2 & 1.92e-2          &  8 \\
 4096 &  64 & No  & No  & 1.89e-4 & 1.92e-2 & 1.92e-2          &  8 \\
\rowcolor{gray!15}
 4096 & 128 & Yes & Yes & 8.86e-4 & 1.28e+0 & \textbf{8.86e-4} & 16 \\
 8192 &  64 & No  & No  & 1.83e-4 & 1.75e-2 & 1.75e-2          &  8 \\
 8192 & 128 & Yes & No  & 1.72e-4 & 1.68e-2 & 1.68e-2          &  8 \\
\bottomrule
\end{tabular}
\end{table*}

\begin{figure*}[tbp]
  \centering
  \includegraphics[width=\textwidth]{accuracy_vs_baselines.png}
  \caption{Log-scale relative RMSE across all 12 workloads. Evolved
    policy (green) tracks FP16 (gray) exactly. Uniform INT8 (orange)
    diverges by orders of magnitude on outlier workloads---RMSE $>$ 1.0
    means the output is noise.}
  \label{fig:accuracy}
\end{figure*}

\subsection{Entropy as the Key Discriminator}

The entropy distribution across all 19{,}488 evaluated blocks is strongly
bimodal (Figure~\ref{fig:entropy}): 13{,}840 blocks (71.0\%) cluster at
entropy $\approx$4.3 (uniform attention, safe for INT8); 5{,}648 blocks
(29.0\%) cluster at entropy 0.5--1.0 (peaked attention from outlier
workloads, assigned FP16). The gap between 1.5 and 3.5 is nearly empty,
making the threshold robust to its exact value.

\begin{figure*}[tbp]
  \centering
  \includegraphics[width=\textwidth]{entropy_distribution.png}
  \caption{Left: Bimodal block entropy histogram with the 2.0 threshold
    marked. Two clusters are clearly separated by a wide gap.
    Right: Precision decisions by entropy---clean separation at the
    threshold. INT8 blocks (blue dots) cluster at high entropy; FP16
    blocks (red crosses) at low entropy.}
  \label{fig:entropy}
\end{figure*}

\begin{figure*}[tbp]
  \centering
  \includegraphics[width=\textwidth]{precision_heatmaps.png}
  \caption{Per-block precision maps across all 12 workloads (blue=INT8,
    red=FP16). Each sub-plot is block\_row~$\times$~block\_col. Outlier
    workloads are uniformly FP16; non-outlier workloads are uniformly
    INT8. No mixed blocks appear within any single workload.}
  \label{fig:heatmaps}
\end{figure*}

\subsection{Sequence Length Invariance}

The policy's decisions are independent of sequence length
(Figure~\ref{fig:scaling}). Across $N = 512$ to $N = 8{,}192$,
compression is constant within each workload family, contradicting the
hypothesis that longer sequences have proportionally more compressible
blocks. The outlier/entropy signal dominates positional effects entirely.

\begin{figure}[tbp]
  \centering
  \includegraphics[width=\columnwidth]{seq_length_scaling.png}
  \caption{Compression vs.\ sequence length. Both lines are flat:
    the precision decision is length-invariant.}
  \label{fig:scaling}
\end{figure}

%------------------------------------------------------------------
\subsection{Comparison to FlashAttention-2}
\label{sec:fa2_comparison}

We benchmark the evolved policy against PyTorch SDPA dispatching to
FlashAttention-2 (FA-2) on an NVIDIA RTX A4000 (Ampere, sm\_86).
Table~\ref{tab:fa2_comparison} reports per-workload relative RMSE
against a FP32 tiled reference, alongside the direct output RMSE
between policy and FA-2 (Match~$<$~5\%).

All 12 workloads pass the match criterion (12/12). On the five
outlier workloads, the policy achieves \emph{lower} RMSE than FA-2:
$0.37$--$0.73\times$ FA-2 error. This advantage arises because our
policy accumulates in FP32 while FA-2 runs end-to-end in FP16---an
effect that is most pronounced when outlier scores amplify FP16
rounding. On the seven non-outlier workloads the policy uses INT8
compression (8 bits, $2\times$ savings) and its RMSE rises to
$\approx$0.018, still within the match threshold. The KV-cache
achieves a uniform $2\times$ compression ($50\%$ memory reduction)
across all workloads and all sequence lengths from 512 to 8{,}192.

\begin{table*}[tbp]
\centering
\caption{Evolved policy vs.\ FlashAttention-2 (PyTorch SDPA,
  RTX A4000). RMSE vs.\ FP32 reference.
  \textbf{Bold}: policy outperforms FA-2.
  All 12/12 workloads: direct output RMSE vs.\ FA-2 $<$ 5\%.}
\label{tab:fa2_comparison}
\small
\begin{tabular}{ccrrrr}
\toprule
$N$ & $d$ & Outliers & FA-2 RMSE & Policy RMSE & Bits \\
\midrule
  512 &  64 & No  & 4.6e-4 & 1.8e-2 &  8 \\
\rowcolor{gray!15}
  512 & 128 & Yes & 1.2e-3 & \textbf{5.7e-4} & 16 \\
\rowcolor{gray!15}
 2048 &  64 & Yes & 2.3e-3 & \textbf{1.0e-3} & 16 \\
 2048 & 128 & No  & 4.2e-4 & 1.8e-2 &  8 \\
 8192 &  64 & No  & 4.3e-4 & 1.6e-2 &  8 \\
\rowcolor{gray!15}
 8192 & 128 & Yes & 2.3e-3 & \textbf{1.1e-3} & 16 \\
\rowcolor{gray!15}
 2048 &  64 & Yes & 2.0e-3 & \textbf{1.1e-3} & 16 \\
 2048 & 128 & No  & 4.6e-4 & 1.9e-2 &  8 \\
 4096 &  64 & No  & 4.7e-4 & 1.9e-2 &  8 \\
\rowcolor{gray!15}
 4096 & 128 & Yes & 1.6e-3 & \textbf{8.9e-4} & 16 \\
 8192 &  64 & No  & 4.7e-4 & 1.8e-2 &  8 \\
 8192 & 128 & No  & 4.3e-4 & 1.7e-2 &  8 \\
\midrule
Average &  &  & 1.0e-3 & 1.1e-2 & 11.3 \\
\bottomrule
\end{tabular}

\end{table*}

%------------------------------------------------------------------
\section{KV-Cache Application}

We implement an entropy-guided KV-cache quantization module with four
hardware-mapped components:

\begin{enumerate}
  \item \textbf{Precision Controller.} Entropy + outlier $\rightarrow$
    1-bit precision select. Hardware: comparator + AND gate.
  \item \textbf{Quantizer.} Symmetric INT8 (max-abs scan,
    multiply-round-clamp) or FP16 passthrough. Hardware: multiplier
    pipeline with precision-select mux.
  \item \textbf{Cache SRAM.} Fixed header:
    \texttt{[1-bit tag][16-bit scale][N$\times$D payload]}.
  \item \textbf{Dequantizer.} Reads tag, applies inverse. Hardware:
    single multiplier pipeline.
\end{enumerate}

Benchmarks across all 12 workloads confirm a uniform $2\times$
compression ratio (50\% memory reduction) at 8 bits/element, scaling
flat from $N=512$ to $N=8{,}192$. This is validated by 45 passing unit
tests covering correctness (MAE $<$0.02), autoregressive append, and
multi-layer storage.

%------------------------------------------------------------------
\section{Hardware Design Implications}

\subsection{Why Simple Beats Learned}

The search had full freedom to evolve arbitrary logic over all 14
features. It converged on two comparisons---evidence that the problem is
genuinely simple. A minimal MLP controller would require multipliers,
weight SRAM, and activation logic on the critical path. The threshold
comparator requires one fixed-point comparator, one flag, and one AND
gate.

\subsection{Programmable Threshold Registers}

We propose storing per-layer thresholds in a small register file:

\smallskip
\noindent\texttt{precision = (has\_outlier \&\& entropy < THRESH[l]) ? FP16 : INT8;}
\smallskip

For a 32-layer transformer this costs 64 bytes. Programmed once at model
load time from offline profiling, this provides per-layer adaptation with
no additional logic---capturing $\sim$95\% of the benefit of a learned
controller at $\sim$0\% of the silicon cost.

\subsection{Entropy Approximation in Hardware}

Entropy accumulates from the log-sum-exp statistics already computed
during FlashAttention's online softmax pass---no extra memory traffic.
A coarser approximation from the max-to-mean ratio of pre-softmax scores
requires only a comparator and divider.

%------------------------------------------------------------------
\section{Limitations}

(1)~\textbf{Synthetic data only}---real LLM activations may shift the
entropy distribution. (2)~\textbf{Workload-level granularity}---no
fine-grained per-block decisions within outlier sequences.
(3)~\textbf{Fixed outlier threshold}---not co-evolved, may not transfer.
(4)~\textbf{No throughput measurement}---accuracy and compression only.
(5)~\textbf{Dense attention only}---grouped-query, sliding window, and
sparse patterns untested.

%------------------------------------------------------------------
\section{Future Work}

(1)~Real-activation validation on Llama, Mistral, Qwen.
(2)~Per-layer entropy profiling for static compile-time assignment.
(3)~RTL prototype: SystemVerilog precision controller with area/power
measurement.
(4)~Co-evolution of outlier and entropy thresholds.
(5)~Triton/CUDA kernel integration with compile-time precision
specialization.

%------------------------------------------------------------------
\section{Conclusion}

Evolutionary search over per-block statistics converges on a simple,
hardware-friendly precision policy for mixed-precision attention. Softmax
entropy is a reliable discriminator between blocks tolerating INT8 and
those requiring FP16. The bimodal distribution provides a wide decision
margin, making the threshold robust. Two comparisons and an AND gate are
sufficient for direct ASIC implementation---no learned controller needed.

%------------------------------------------------------------------
\bibliographystyle{plainnat}
\bibliography{references}

\end{document}
